\latinkeywords{Implicit Sentiment Analysis, Large Language Models, Chain-of-Thought, Multi-Stage Reasoning, Source Selection}

\en-abstract{
Detecting sentiment in texts without explicit evaluative words requires contextual understanding, commonsense knowledge, and inference about the relationship between a described event and its target. This thesis presents a prompt-based, multi-source system for implicit sentiment analysis. The system combines a language model's direct prediction with multi-stage reasoning inspired by THOR, structured error diagnosis, and a source-selection policy. The final selection policy is tuned using the training data and, without access to gold test labels, selects one of the direct, reasoning, or diagnostic sources. The main evaluation uses the SCAPT-labelled data derived from SemEval 2014 in the laptop and restaurant domains. The main dataset contains 2,188 examples, with 1,746 training examples and 442 test examples. The primary experiments use Qwen3 8B through Ollama. On the test set, the final system increases accuracy from 0.678733 to 0.723982 and Macro-F1 from 0.674075 to 0.719119. A supplementary comparison with Gemini 2.5 Flash is also conducted on a balanced subset of 240 examples. For the model, data, and experimental setting examined in this thesis, the results indicate that controlled source selection is more effective than relying on one fixed prediction path.
}

\newpage
\thispagestyle{empty}
\begin{latin}
\section*{\LARGE\centering Abstract}
\een-abstract
\vspace*{.5cm}
{\large\textbf{Key Words:}}\par
\vspace*{.5cm}
\elatinkeywords
\end{latin}
